%% kapa: exam=1 round=22 year=2011 period=1 form=A subject=economics subject_ko=경제원론 seq=37 number=77
%% source: Q-Net 감정평가사 기출문제 (cst00309) artlSeq=5036790 file=extracted/1008/5036790/2011년 제22회 감정평가사 1차시험 기출문제/1교시.hwp
%% conversion: kordoc 4.24.51 -> tools/gen1.py (auto); review=proofread-sonnet; flags=none
\begin{question}{77}
어빙 피셔(Irving Fisher)의 2기간 최적 소비선택 모형에서 도출되는 결론으로 옳은 것만을 모두 고른 것은? (단, 기간별로 소비되는 재화는 모두 정상재, 차입제약은 없고, 각 기간의 소비는 모두 0보다 큼)
\begin{framedtext}
ㄱ. 제1기의 소득증가는 제1기의 소비를 증가시킨다. \par
ㄴ. 제2기의 소득증가는 제2기의 소비를 감소시킨다. \par
ㄷ. 실질이자율이 상승하면 제2기의 소비는 증가한다. \par
ㄹ. 제2기의 소득증가는 제1기의 소비를 감소시킨다.
\end{framedtext}
\begin{choices}
  \item ㄱ, ㄴ
  \item ㄴ, ㄷ
  \item ㄱ, ㄷ
  \item ㄴ, ㄹ
  \item ㄷ, ㄹ
\end{choices}
\end{question}
\answerall
\answernote{정답표에 1,2,3,4,5 모두 정답으로 표기(전원 정답)}
